\documentclass[12pt,]{article}
\usepackage{lmodern}
\usepackage{amssymb,amsmath}
\usepackage{ifxetex,ifluatex}
\usepackage{fixltx2e} % provides \textsubscript
\ifnum 0\ifxetex 1\fi\ifluatex 1\fi=0 % if pdftex
  \usepackage[T1]{fontenc}
  \usepackage[utf8]{inputenc}
\else % if luatex or xelatex
  \ifxetex
    \usepackage{mathspec}
  \else
    \usepackage{fontspec}
  \fi
  \defaultfontfeatures{Ligatures=TeX,Scale=MatchLowercase}
\fi
% use upquote if available, for straight quotes in verbatim environments
\IfFileExists{upquote.sty}{\usepackage{upquote}}{}
% use microtype if available
\IfFileExists{microtype.sty}{%
\usepackage{microtype}
\UseMicrotypeSet[protrusion]{basicmath} % disable protrusion for tt fonts
}{}
\usepackage[margin=1in]{geometry}
\usepackage{hyperref}
\hypersetup{unicode=true,
            pdftitle={In Praise of QR},
            pdfauthor={Jan de Leeuw},
            pdfborder={0 0 0},
            breaklinks=true}
\urlstyle{same}  % don't use monospace font for urls
\usepackage{color}
\usepackage{fancyvrb}
\newcommand{\VerbBar}{|}
\newcommand{\VERB}{\Verb[commandchars=\\\{\}]}
\DefineVerbatimEnvironment{Highlighting}{Verbatim}{commandchars=\\\{\}}
% Add ',fontsize=\small' for more characters per line
\usepackage{framed}
\definecolor{shadecolor}{RGB}{248,248,248}
\newenvironment{Shaded}{\begin{snugshade}}{\end{snugshade}}
\newcommand{\KeywordTok}[1]{\textcolor[rgb]{0.13,0.29,0.53}{\textbf{#1}}}
\newcommand{\DataTypeTok}[1]{\textcolor[rgb]{0.13,0.29,0.53}{#1}}
\newcommand{\DecValTok}[1]{\textcolor[rgb]{0.00,0.00,0.81}{#1}}
\newcommand{\BaseNTok}[1]{\textcolor[rgb]{0.00,0.00,0.81}{#1}}
\newcommand{\FloatTok}[1]{\textcolor[rgb]{0.00,0.00,0.81}{#1}}
\newcommand{\ConstantTok}[1]{\textcolor[rgb]{0.00,0.00,0.00}{#1}}
\newcommand{\CharTok}[1]{\textcolor[rgb]{0.31,0.60,0.02}{#1}}
\newcommand{\SpecialCharTok}[1]{\textcolor[rgb]{0.00,0.00,0.00}{#1}}
\newcommand{\StringTok}[1]{\textcolor[rgb]{0.31,0.60,0.02}{#1}}
\newcommand{\VerbatimStringTok}[1]{\textcolor[rgb]{0.31,0.60,0.02}{#1}}
\newcommand{\SpecialStringTok}[1]{\textcolor[rgb]{0.31,0.60,0.02}{#1}}
\newcommand{\ImportTok}[1]{#1}
\newcommand{\CommentTok}[1]{\textcolor[rgb]{0.56,0.35,0.01}{\textit{#1}}}
\newcommand{\DocumentationTok}[1]{\textcolor[rgb]{0.56,0.35,0.01}{\textbf{\textit{#1}}}}
\newcommand{\AnnotationTok}[1]{\textcolor[rgb]{0.56,0.35,0.01}{\textbf{\textit{#1}}}}
\newcommand{\CommentVarTok}[1]{\textcolor[rgb]{0.56,0.35,0.01}{\textbf{\textit{#1}}}}
\newcommand{\OtherTok}[1]{\textcolor[rgb]{0.56,0.35,0.01}{#1}}
\newcommand{\FunctionTok}[1]{\textcolor[rgb]{0.00,0.00,0.00}{#1}}
\newcommand{\VariableTok}[1]{\textcolor[rgb]{0.00,0.00,0.00}{#1}}
\newcommand{\ControlFlowTok}[1]{\textcolor[rgb]{0.13,0.29,0.53}{\textbf{#1}}}
\newcommand{\OperatorTok}[1]{\textcolor[rgb]{0.81,0.36,0.00}{\textbf{#1}}}
\newcommand{\BuiltInTok}[1]{#1}
\newcommand{\ExtensionTok}[1]{#1}
\newcommand{\PreprocessorTok}[1]{\textcolor[rgb]{0.56,0.35,0.01}{\textit{#1}}}
\newcommand{\AttributeTok}[1]{\textcolor[rgb]{0.77,0.63,0.00}{#1}}
\newcommand{\RegionMarkerTok}[1]{#1}
\newcommand{\InformationTok}[1]{\textcolor[rgb]{0.56,0.35,0.01}{\textbf{\textit{#1}}}}
\newcommand{\WarningTok}[1]{\textcolor[rgb]{0.56,0.35,0.01}{\textbf{\textit{#1}}}}
\newcommand{\AlertTok}[1]{\textcolor[rgb]{0.94,0.16,0.16}{#1}}
\newcommand{\ErrorTok}[1]{\textcolor[rgb]{0.64,0.00,0.00}{\textbf{#1}}}
\newcommand{\NormalTok}[1]{#1}
\usepackage{graphicx,grffile}
\makeatletter
\def\maxwidth{\ifdim\Gin@nat@width>\linewidth\linewidth\else\Gin@nat@width\fi}
\def\maxheight{\ifdim\Gin@nat@height>\textheight\textheight\else\Gin@nat@height\fi}
\makeatother
% Scale images if necessary, so that they will not overflow the page
% margins by default, and it is still possible to overwrite the defaults
% using explicit options in \includegraphics[width, height, ...]{}
\setkeys{Gin}{width=\maxwidth,height=\maxheight,keepaspectratio}
\IfFileExists{parskip.sty}{%
\usepackage{parskip}
}{% else
\setlength{\parindent}{0pt}
\setlength{\parskip}{6pt plus 2pt minus 1pt}
}
\setlength{\emergencystretch}{3em}  % prevent overfull lines
\providecommand{\tightlist}{%
  \setlength{\itemsep}{0pt}\setlength{\parskip}{0pt}}
\setcounter{secnumdepth}{5}
% Redefines (sub)paragraphs to behave more like sections
\ifx\paragraph\undefined\else
\let\oldparagraph\paragraph
\renewcommand{\paragraph}[1]{\oldparagraph{#1}\mbox{}}
\fi
\ifx\subparagraph\undefined\else
\let\oldsubparagraph\subparagraph
\renewcommand{\subparagraph}[1]{\oldsubparagraph{#1}\mbox{}}
\fi

%%% Use protect on footnotes to avoid problems with footnotes in titles
\let\rmarkdownfootnote\footnote%
\def\footnote{\protect\rmarkdownfootnote}

%%% Change title format to be more compact
\usepackage{titling}

% Create subtitle command for use in maketitle
\providecommand{\subtitle}[1]{
  \posttitle{
    \begin{center}\large#1\end{center}
    }
}

\setlength{\droptitle}{-2em}

  \title{In Praise of QR}
    \pretitle{\vspace{\droptitle}\centering\huge}
  \posttitle{\par}
    \author{Jan de Leeuw}
    \preauthor{\centering\large\emph}
  \postauthor{\par}
      \predate{\centering\large\emph}
  \postdate{\par}
    \date{Version 003, March 21, 2016}


\begin{document}
\maketitle

{
\setcounter{tocdepth}{3}
\tableofcontents
}
Note: This is a working paper which will be expanded/updated frequently.
Suggestions for improvement are welcome. The directory
\href{http://deleeuwpdx.net/pubfolders/matrix}{deleeuwpdx.net/pubfolders/matrix}
has a pdf copy of this article, the complete Rmd file with all code
chunks, the bib file, and the R and C source code. We thank Bill
Venables for some corrections in the handling of pivots in version 001.

\section{Problem}\label{problem}

Consider the linear system \(Xb=y\), with \(X\) an \(n\times m\) matrix.
The function \texttt{solve()} in \texttt{R} handles the case in which
\(X\) is square and non-singular, while \texttt{solve.qr()} handles
rectangular systems, providing the least squares solution if
appropriate. Of course a solution of \(Xb=y\) may not exist, and if it
exists it may not be unique. A \emph{least squares solution}, which
minimizes the sum of squares of the residuals \(y-Xb\), always exists,
but again it may not be unique. If \(Xb=y\) is solvable, then of course
the solutions are the least squares solutions, and the residual sum of
squares is zero. In this note we provide software in \texttt{R} and
\texttt{C} for computing a basis for the affine space of all solutions
and/or all least squares solutions to the system (including the case
where the solution set is empty). In addition we provide software to
compute the Moore-Penrose inverse and to compute an orthonormal bases
for the null-space of \(X\).

\section{The QR decomposition}\label{the-qr-decomposition}

\subsection{Theory}\label{theory}

At the core of our treatment of linear systems is the QR decomposition.
We state it in a rank revealing form. This is theorem 1.3.4 in Björck
(1996), but the proof can be found in most numerical linear algebra
books.

\textbf{Theorem 1: {[}QR{]}} If the \(n\times m\) matrix \(X\) has rank
\(r\) then there exist

\begin{itemize}
\tightlist
\item
  An \(n\times r\) matrix \(Q\) with \(Q'Q=I\),
\item
  An \(r\times r\) upper-triangular matrix \(T\) with positive diagonal
  elements,
\item
  An \(r\times(m-r)\) matrix \(S\),
\item
  An \(m\times m\) permutation matrix \(P\),
\end{itemize}

such that

\begin{equation}
XP=QR,\label{E:qr}
\end{equation}

with \(R=(T\mid S)\).

The permutation matrix \(P\) is used to makes sure that the leading
\(n\times r\) submatrix of \(XP\) is of rank \(r\). The theorem says the
QR decomposition is a \emph{full rank decomposition}, because both \(Q\)
and \(R:=(T \mid S)\) are of rank \(r=\mathbf{rank}(X)\).

We have written the QR decomposition the form in which it is computed by
our code. Note that alternatively we can write it in a more verbose way
as

\begin{equation}
X=\begin{bmatrix}Q&\mid&Q_\perp\end{bmatrix}\begin{bmatrix}T&\mid&S\\--&&--\\0_{(n-r)\times r}&\mid&0_{(n-r)\times(m-r)}\end{bmatrix}\ P',\label{E:qr2}
\end{equation}

with \(Q_\perp\) an \(n\times(n-r)\) orthonormal basis for the
complement of the column space of \(X\) (i.e.~the \emph{left null space}
of \(X\), the set of all \(y\) such that \(y'X=0\)).

\subsection{Computation}\label{computation}

In De Leeuw (2015) we computed the QR decomposition using the
\emph{modified Gram-Schmidt (MGS) algorithm}, which is described, for
example, in section 2.4.2 of Björck (1996). The software uses the good
old \texttt{.C()} interface in \texttt{R}. It is \emph{exceedingly
simple} because it can only handle the case where \(n\geq m=r\), i.e.
\(A\) has to be \emph{tall} and \emph{non-singular}. This is fine for
some applications, but useless for what we have in mind in this paper.

We have consequently written a new MGS routine, which computes the
rank-revealing QR decomposition from theorem 1. The \texttt{R} function
\texttt{gsRC()} sets up the memory, passes the arguments using the
\texttt{.C()} interface to the \texttt{C} function \texttt{gsC()}, and
then extracts the output from the returned list. The function
\texttt{gsC()} can be easily incorporated into both \texttt{R} and
\texttt{C} projects. It passes by reference, does not use input and
output routines, does not allocate or free memory, does not call
\texttt{R} functions, does not use \texttt{C} arrays, and does not
return a value. All I/O and memory management must be handled by the
\texttt{C} or \texttt{R} calling program. This implies we also pass
pointers to memory for the results. All \texttt{R} matrices are passed
as vectors and the \texttt{C} routine takes into account they are in
column-major order with index starting at one instead of zero. A simple
\texttt{C} macro \texttt{MINDEX()} handles all array indexing.

The rank reducing QR decomposition requires \emph{column-pivoting} to
find an appropriate permutation matrix \(P\). In MGS we handle the
columns of \(X\) in order, so each column becomes the \emph{working
column} in turn. We start a counter for the \emph{rank} at \(m\) and a
counter for the \emph{working number} at 1. The working column is
projected on the orthogonal complement of the space spanned by the
previous working columns. If the working column is linearly dependent on
the previous columns, i.e.~if the projection is zero, then we exchange
it with a column with a higher column number and decrease the rank by
one. The working column number stays the same. If the working column is
independent of the previous ones we normalize the projection and add it
to \(Q\). The rank counter stays the same, the working number increases
by one. In the MGS algorithm \(A\) is transformed in place, no extra
storage for \(Q\) is needed.

Since the pivoting, and the resulting column permutations, may be
confusing, we show the results of our function \texttt{gsRC()} on a
singular example in some detail. The construction and manipulation of
the permutation matrix \(P\) from the pivot sequence are also
illustrated.

\begin{Shaded}
\begin{Highlighting}[]
\KeywordTok{print}\NormalTok{ (x <-}\StringTok{ }\KeywordTok{matrix}\NormalTok{(}\KeywordTok{c}\NormalTok{(}\DecValTok{1}\NormalTok{,}\DecValTok{1}\NormalTok{,}\DecValTok{1}\NormalTok{,}\DecValTok{1}\NormalTok{,}\DecValTok{1}\NormalTok{,}\OperatorTok{-}\DecValTok{1}\NormalTok{,}\DecValTok{1}\NormalTok{,}\OperatorTok{-}\DecValTok{1}\NormalTok{,}\DecValTok{2}\NormalTok{,}\DecValTok{0}\NormalTok{,}\DecValTok{2}\NormalTok{,}\DecValTok{0}\NormalTok{,}\DecValTok{1}\NormalTok{,}\OperatorTok{-}\DecValTok{1}\NormalTok{,}\OperatorTok{-}\DecValTok{1}\NormalTok{,}\DecValTok{1}\NormalTok{,}\DecValTok{0}\NormalTok{,}\DecValTok{2}\NormalTok{,}\DecValTok{0}\NormalTok{,}\DecValTok{2}\NormalTok{), }\DecValTok{4}\NormalTok{, }\DecValTok{5}\NormalTok{))}
\end{Highlighting}
\end{Shaded}

\begin{verbatim}
##      [,1] [,2] [,3] [,4] [,5]
## [1,]    1    1    2    1    0
## [2,]    1   -1    0   -1    2
## [3,]    1    1    2   -1    0
## [4,]    1   -1    0    1    2
\end{verbatim}

\begin{Shaded}
\begin{Highlighting}[]
\KeywordTok{print}\NormalTok{ (h <-}\StringTok{ }\KeywordTok{gsRC}\NormalTok{(x))}
\end{Highlighting}
\end{Shaded}

\begin{verbatim}
## $x
##      [,1] [,2] [,3]
## [1,]  0.5  0.5  0.5
## [2,]  0.5 -0.5 -0.5
## [3,]  0.5  0.5 -0.5
## [4,]  0.5 -0.5  0.5
## 
## $r
##      [,1] [,2] [,3] [,4] [,5]
## [1,]    2    0    0    2    2
## [2,]    0    2    0   -2    2
## [3,]    0    0    2    0    0
## 
## $rank
## [1] 3
## 
## $pivot
## [1] 1 2 4 5 3
\end{verbatim}

\begin{Shaded}
\begin{Highlighting}[]
\KeywordTok{print}\NormalTok{ (p <-}\StringTok{ }\KeywordTok{ifelse}\NormalTok{(}\KeywordTok{outer}\NormalTok{(}\DecValTok{1}\OperatorTok{:}\DecValTok{5}\NormalTok{, h}\OperatorTok{$}\NormalTok{pivot, }\StringTok{"=="}\NormalTok{), }\DecValTok{1}\NormalTok{,}\DecValTok{0}\NormalTok{))}
\end{Highlighting}
\end{Shaded}

\begin{verbatim}
##      [,1] [,2] [,3] [,4] [,5]
## [1,]    1    0    0    0    0
## [2,]    0    1    0    0    0
## [3,]    0    0    0    0    1
## [4,]    0    0    1    0    0
## [5,]    0    0    0    1    0
\end{verbatim}

\begin{Shaded}
\begin{Highlighting}[]
\NormalTok{x }\OperatorTok{%*%}\StringTok{ }\NormalTok{p}
\end{Highlighting}
\end{Shaded}

\begin{verbatim}
##      [,1] [,2] [,3] [,4] [,5]
## [1,]    1    1    1    0    2
## [2,]    1   -1   -1    2    0
## [3,]    1    1   -1    0    2
## [4,]    1   -1    1    2    0
\end{verbatim}

\begin{Shaded}
\begin{Highlighting}[]
\NormalTok{h}\OperatorTok{$}\NormalTok{x }\OperatorTok{%*%}\StringTok{ }\NormalTok{h}\OperatorTok{$}\NormalTok{r}
\end{Highlighting}
\end{Shaded}

\begin{verbatim}
##      [,1] [,2] [,3] [,4] [,5]
## [1,]    1    1    1    0    2
## [2,]    1   -1   -1    2    0
## [3,]    1    1   -1    0    2
## [4,]    1   -1    1    2    0
\end{verbatim}

\begin{Shaded}
\begin{Highlighting}[]
\NormalTok{x[, h}\OperatorTok{$}\NormalTok{pivot]}
\end{Highlighting}
\end{Shaded}

\begin{verbatim}
##      [,1] [,2] [,3] [,4] [,5]
## [1,]    1    1    1    0    2
## [2,]    1   -1   -1    2    0
## [3,]    1    1   -1    0    2
## [4,]    1   -1    1    2    0
\end{verbatim}

\begin{Shaded}
\begin{Highlighting}[]
\NormalTok{(h}\OperatorTok{$}\NormalTok{x }\OperatorTok{%*%}\StringTok{ }\NormalTok{h}\OperatorTok{$}\NormalTok{r)[, }\KeywordTok{order}\NormalTok{(h}\OperatorTok{$}\NormalTok{pivot)]}
\end{Highlighting}
\end{Shaded}

\begin{verbatim}
##      [,1] [,2] [,3] [,4] [,5]
## [1,]    1    1    2    1    0
## [2,]    1   -1    0   -1    2
## [3,]    1    1    2   -1    0
## [4,]    1   -1    0    1    2
\end{verbatim}

\begin{Shaded}
\begin{Highlighting}[]
\KeywordTok{tcrossprod}\NormalTok{ (h}\OperatorTok{$}\NormalTok{x }\OperatorTok{%*%}\StringTok{ }\NormalTok{h}\OperatorTok{$}\NormalTok{r, p)}
\end{Highlighting}
\end{Shaded}

\begin{verbatim}
##      [,1] [,2] [,3] [,4] [,5]
## [1,]    1    1    2    1    0
## [2,]    1   -1    0   -1    2
## [3,]    1    1    2   -1    0
## [4,]    1   -1    0    1    2
\end{verbatim}

\subsection{Examples}\label{examples}

\subsubsection{Regular}\label{regular}

\begin{Shaded}
\begin{Highlighting}[]
\KeywordTok{set.seed}\NormalTok{(}\DecValTok{12345}\NormalTok{)}
\NormalTok{x <-}\StringTok{ }\KeywordTok{matrix}\NormalTok{ (}\KeywordTok{rnorm}\NormalTok{(}\DecValTok{20}\NormalTok{), }\DecValTok{5}\NormalTok{, }\DecValTok{4}\NormalTok{)}
\KeywordTok{print}\NormalTok{(}\KeywordTok{lapply}\NormalTok{ (}\KeywordTok{gsRC}\NormalTok{ (x), zapsmall), }\DataTypeTok{digits =} \DecValTok{4}\NormalTok{)}
\end{Highlighting}
\end{Shaded}

\begin{verbatim}
## $x
##          [,1]    [,2]    [,3]     [,4]
## [1,]  0.48949 -0.7027  0.2543 -0.04908
## [2,]  0.59310  0.5679  0.5599  0.08871
## [3,] -0.09138 -0.1772  0.3475 -0.79043
## [4,] -0.37912 -0.3036  0.5817  0.56200
## [5,]  0.50651 -0.2452 -0.4034  0.22161
## 
## $r
##       [,1]    [,2]   [,3]    [,4]
## [1,] 1.196 -0.8488 0.4097 -0.3691
## [2,] 0.000  1.9959 1.0742 -1.4321
## [3,] 0.000  0.0000 1.7221  0.1276
## [4,] 0.000  0.0000 0.0000  0.8394
## 
## $rank
## [1] 4
## 
## $pivot
## [1] 1 2 3 4
\end{verbatim}

\subsubsection{Singular}\label{singular}

\begin{Shaded}
\begin{Highlighting}[]
\KeywordTok{set.seed}\NormalTok{(}\DecValTok{12345}\NormalTok{)}
\NormalTok{x <-}\StringTok{ }\KeywordTok{matrix}\NormalTok{ (}\KeywordTok{rnorm}\NormalTok{(}\DecValTok{20}\NormalTok{), }\DecValTok{5}\NormalTok{, }\DecValTok{4}\NormalTok{)}
\NormalTok{x[,}\DecValTok{3}\NormalTok{] <-}\StringTok{ }\NormalTok{x[,}\DecValTok{1}\NormalTok{] }\OperatorTok{+}\StringTok{ }\NormalTok{x[,}\DecValTok{2}\NormalTok{]}
\KeywordTok{print}\NormalTok{(}\KeywordTok{lapply}\NormalTok{ (}\KeywordTok{gsRC}\NormalTok{ (x), zapsmall), }\DataTypeTok{digits =} \DecValTok{4}\NormalTok{)}
\end{Highlighting}
\end{Shaded}

\begin{verbatim}
## $x
##          [,1]    [,2]     [,3]
## [1,]  0.48949 -0.7027 -0.01029
## [2,]  0.59310  0.5679  0.17187
## [3,] -0.09138 -0.1772 -0.72921
## [4,] -0.37912 -0.3036  0.64304
## [5,]  0.50651 -0.2452  0.15846
## 
## $r
##       [,1]    [,2]    [,3]   [,4]
## [1,] 1.196 -0.8488 -0.3691 0.3474
## [2,] 0.000  1.9959 -1.4321 1.9959
## [3,] 0.000  0.0000  0.8490 0.0000
## 
## $rank
## [1] 3
## 
## $pivot
## [1] 1 2 4 3
\end{verbatim}

\subsubsection{Broad}\label{broad}

\begin{Shaded}
\begin{Highlighting}[]
\KeywordTok{set.seed}\NormalTok{(}\DecValTok{12345}\NormalTok{)}
\NormalTok{x <-}\StringTok{ }\KeywordTok{matrix}\NormalTok{ (}\KeywordTok{rnorm}\NormalTok{(}\DecValTok{20}\NormalTok{), }\DecValTok{5}\NormalTok{, }\DecValTok{4}\NormalTok{)}
\KeywordTok{print}\NormalTok{(}\KeywordTok{lapply}\NormalTok{ (}\KeywordTok{gsRC}\NormalTok{ (}\KeywordTok{t}\NormalTok{(x)), zapsmall), }\DataTypeTok{digits =} \DecValTok{4}\NormalTok{)}
\end{Highlighting}
\end{Shaded}

\begin{verbatim}
## $x
##          [,1]     [,2]    [,3]     [,4]
## [1,]  0.28143  0.44828 -0.6526 -0.54221
## [2,] -0.87379 -0.03325 -0.4763  0.09223
## [3,] -0.05587  0.85010  0.1408  0.50436
## [4,]  0.39264 -0.27435 -0.5722  0.66567
## 
## $r
##       [,1]    [,2]    [,3]     [,4]    [,5]
## [1,] 2.081 -0.8005 0.05967  0.53164  1.1330
## [2,] 0.000  2.0852 0.36622 -0.05908 -0.4178
## [3,] 0.000  0.0000 0.44481 -0.13676 -0.2341
## [4,] 0.000  0.0000 0.00000  1.22809 -0.5930
## 
## $rank
## [1] 4
## 
## $pivot
## [1] 1 2 3 4 5
\end{verbatim}

\subsubsection{Weird}\label{weird}

\begin{Shaded}
\begin{Highlighting}[]
\NormalTok{x <-}\StringTok{ }\KeywordTok{matrix}\NormalTok{(}\DecValTok{1}\NormalTok{, }\DecValTok{1}\NormalTok{, }\DecValTok{6}\NormalTok{)}
\KeywordTok{print}\NormalTok{(}\KeywordTok{lapply}\NormalTok{ (}\KeywordTok{gsRC}\NormalTok{ (x), zapsmall), }\DataTypeTok{digits =} \DecValTok{4}\NormalTok{)}
\end{Highlighting}
\end{Shaded}

\begin{verbatim}
## $x
##      [,1]
## [1,]    1
## 
## $r
##      [,1] [,2] [,3] [,4] [,5] [,6]
## [1,]    1    1    1    1    1    1
## 
## $rank
## [1] 1
## 
## $pivot
## [1] 1 3 4 5 6 2
\end{verbatim}

\section{Least Squares}\label{least-squares}

\subsection{Theory}\label{theory-1}

Least squares solutions are minimizers of the sum of squares, which we
write as \(\mathbf{SSQ}\ (y-Xb)\). For solvable systems \(Xb=y\) the set
of least squares solutions and the set of solutions to the system are
obviously identical.

\textbf{Theorem 2: {[}LS{]}} Suppose \(X\) has QR decomposition \(QRP'\)
of theorem 1. Then

\begin{itemize}
\tightlist
\item
  the minimum of \(\mathbf{SSQ}\ (y- Xb)\) is \(y'(I-QQ')y\),
\item
  the least squares solutions are all are of the form
\end{itemize}

\begin{equation}
b = P\left\{\begin{bmatrix}T^{-1}Q'y\\0\end{bmatrix}+\begin{bmatrix}-T^{-1}S\\I\end{bmatrix}v\right\}.\label{E:ls}
\end{equation}

\textbf{Proof:} We can write \(y-Xb=Q(Q'y-R\tilde b)+(I-QQ')y\) with
\(\tilde b:=P'b\), and thus \[
\mathbf{SSQ}\ (y-Xb)=y'(I-QQ')y+\mathbf{SSQ}\ (Q'y - T\tilde b_1 - S\tilde b_2).
\] The second term can be always made equal to zero by choosing
\(\tilde b_2\) arbitrarily and setting \(\tilde b_1=T^{-1}(Q'y-Sb_2)\).
\textbf{QED}

The theorem has a simple special case, in which the residual sum of
squares is zero.

\textbf{Corollary 1: {[}Linear{]}} Suppose \(X\) has the QR
decomposition \(QRP'\) of theorem 1. Then

\begin{itemize}
\tightlist
\item
  \(Xb=y\) is solvable if and only if \(QQ'y=y\).
\item
  If \(Xb=y\) is solvable, then all solutions are of the form
\end{itemize}

\begin{equation}
b = P\left\{\begin{bmatrix}T^{-1}Q'y\\0\end{bmatrix}+\begin{bmatrix}-T^{-1}S\\I\end{bmatrix}v\right\}.\label{E:linear}
\end{equation}

\textbf{Proof:} Directly from theorem 2. \textbf{QED}

\subsection{Computation}\label{computation-1}

Our function \texttt{lsRC()} computes least squares solutions by using
the result in theorem 2. It can, of course, also be used to solve
consistent systems of linear equations.

\subsection{Examples}\label{examples-1}

\subsubsection{Regular}\label{regular-1}

\begin{Shaded}
\begin{Highlighting}[]
\KeywordTok{set.seed}\NormalTok{(}\DecValTok{12345}\NormalTok{)}
\NormalTok{x <-}\StringTok{ }\KeywordTok{matrix}\NormalTok{ (}\KeywordTok{rnorm}\NormalTok{(}\DecValTok{20}\NormalTok{), }\DecValTok{5}\NormalTok{, }\DecValTok{4}\NormalTok{)}
\KeywordTok{print}\NormalTok{(}\KeywordTok{lapply}\NormalTok{ (}\KeywordTok{lsRC}\NormalTok{ (x, }\KeywordTok{rep}\NormalTok{(}\DecValTok{1}\NormalTok{, }\DecValTok{5}\NormalTok{)), zapsmall), }\DataTypeTok{digits =} \DecValTok{4}\NormalTok{)}
\end{Highlighting}
\end{Shaded}

\begin{verbatim}
## $solution
## [1]  0.09947 -0.82045  0.77524  0.03908
## 
## $residuals
## [1] -0.49160  0.07219  0.50991  0.36487  0.75564
## 
## $minssq
## [1] 1.211
## 
## $nullspace
##      [,1]
## [1,]    0
## [2,]    0
## [3,]    0
## [4,]    0
## 
## $rank
## [1] 4
## 
## $pivot
## [1] 1 2 3 4
\end{verbatim}

\subsubsection{Singular}\label{singular-1}

\begin{Shaded}
\begin{Highlighting}[]
\KeywordTok{set.seed}\NormalTok{(}\DecValTok{12345}\NormalTok{)}
\NormalTok{x <-}\StringTok{ }\KeywordTok{matrix}\NormalTok{ (}\KeywordTok{rnorm}\NormalTok{(}\DecValTok{20}\NormalTok{), }\DecValTok{5}\NormalTok{, }\DecValTok{4}\NormalTok{)}
\NormalTok{x[,}\DecValTok{3}\NormalTok{] <-}\StringTok{ }\NormalTok{x[,}\DecValTok{1}\NormalTok{] }\OperatorTok{+}\StringTok{ }\NormalTok{x[,}\DecValTok{2}\NormalTok{]}
\KeywordTok{print}\NormalTok{(}\KeywordTok{lapply}\NormalTok{ (}\KeywordTok{lsRC}\NormalTok{ (x, }\KeywordTok{rep}\NormalTok{(}\DecValTok{1}\NormalTok{,}\DecValTok{5}\NormalTok{)), zapsmall), }\DataTypeTok{digits =} \DecValTok{4}\NormalTok{)}
\end{Highlighting}
\end{Shaded}

\begin{verbatim}
## $solution
## [1]  0.8543 -0.2336  0.0000  0.2754
## 
## $residuals
## [1] -0.1500  0.7852  1.1202  1.0124  0.1853
## 
## $minssq
## [1] 2.953
## 
## $nullspace
##      [,1]
## [1,]   -1
## [2,]   -1
## [3,]    1
## [4,]    0
## 
## $rank
## [1] 3
## 
## $pivot
## [1] 1 2 4 3
\end{verbatim}

\subsubsection{Broad}\label{broad-1}

\begin{Shaded}
\begin{Highlighting}[]
\KeywordTok{set.seed}\NormalTok{(}\DecValTok{12345}\NormalTok{)}
\NormalTok{x <-}\StringTok{ }\KeywordTok{matrix}\NormalTok{ (}\KeywordTok{rnorm}\NormalTok{(}\DecValTok{20}\NormalTok{), }\DecValTok{5}\NormalTok{, }\DecValTok{4}\NormalTok{)}
\KeywordTok{print}\NormalTok{(}\KeywordTok{lapply}\NormalTok{ (}\KeywordTok{lsRC}\NormalTok{ (}\KeywordTok{t}\NormalTok{(x), }\KeywordTok{rep}\NormalTok{(}\DecValTok{1}\NormalTok{,}\DecValTok{4}\NormalTok{)), zapsmall), }\DataTypeTok{digits =} \DecValTok{4}\NormalTok{)}
\end{Highlighting}
\end{Shaded}

\begin{verbatim}
## $solution
## [1]  0.2368  1.0762 -3.3275  0.5863  0.0000
## 
## $residuals
## [1] 0.00e+00 0.00e+00 2.22e-16 0.00e+00
## 
## $minssq
## [1] 4.93e-32
## 
## $nullspace
##          [,1]
## [1,] -0.65057
## [2,]  0.09553
## [3,]  0.67480
## [4,]  0.48286
## [5,]  1.00000
## 
## $rank
## [1] 4
## 
## $pivot
## [1] 1 2 3 4 5
\end{verbatim}

\subsubsection{Weird}\label{weird-1}

\begin{Shaded}
\begin{Highlighting}[]
\NormalTok{x <-}\StringTok{ }\KeywordTok{matrix}\NormalTok{(}\DecValTok{1}\NormalTok{, }\DecValTok{1}\NormalTok{, }\DecValTok{6}\NormalTok{)}
\KeywordTok{print}\NormalTok{(}\KeywordTok{lapply}\NormalTok{ (}\KeywordTok{lsRC}\NormalTok{ (x, }\DecValTok{1}\NormalTok{), zapsmall), }\DataTypeTok{digits =} \DecValTok{4}\NormalTok{)}
\end{Highlighting}
\end{Shaded}

\begin{verbatim}
## $solution
## [1] 1 0 0 0 0 0
## 
## $residuals
## [1] 0
## 
## $minssq
## [1] 0
## 
## $nullspace
##      [,1] [,2] [,3] [,4] [,5]
## [1,]   -1   -1   -1   -1   -1
## [2,]    0    0    0    0    1
## [3,]    1    0    0    0    0
## [4,]    0    1    0    0    0
## [5,]    0    0    1    0    0
## [6,]    0    0    0    1    0
## 
## $rank
## [1] 1
## 
## $pivot
## [1] 1 6 2 3 4 5
\end{verbatim}

\section{The Null-space}\label{the-null-space}

\subsection{Theory}\label{theory-2}

Of course \(Xb=0\) always has the solution \(b=0\). The set of all \(z\)
such that \(Xz=0\) is the \emph{right null space} of \(X\), which is
equal to the left null space of the transpose of \(X\). The QR
decomposition allows us to give a more complete description of these
null spaces. Note that \texttt{Null()} in the package MASS (Venables and
Ripley (2002)) computes the right null space, using the singular value
decomposition.

\textbf{Theorem 3: {[}Null{]}} Suppose \(X\) has QR decomposition
\(QRP'\) of theorem 1. Then the columns of

\begin{equation}
P\begin{bmatrix}-T^{-1}S\\I\end{bmatrix}\label{E:null}
\end{equation}

are a basis for the left null-space of \(X\).

\textbf{Proof:} Special case of theorem 2. \textbf{QED}

\subsection{Computation}\label{computation-2}

The basis in theorem 3 is generally not orthonormal. It can easily be
made orthonormal by applying the QR decomposition a second time, this
time to the basis. Our function \texttt{nullRC()} does exactly this.

\subsection{Examples}\label{examples-2}

\subsubsection{Regular}\label{regular-2}

\begin{Shaded}
\begin{Highlighting}[]
\KeywordTok{set.seed}\NormalTok{(}\DecValTok{12345}\NormalTok{)}
\NormalTok{x <-}\StringTok{ }\KeywordTok{matrix}\NormalTok{ (}\KeywordTok{rnorm}\NormalTok{(}\DecValTok{20}\NormalTok{), }\DecValTok{5}\NormalTok{, }\DecValTok{4}\NormalTok{)}
\KeywordTok{print}\NormalTok{(}\KeywordTok{zapsmall}\NormalTok{(}\KeywordTok{nullRC}\NormalTok{ (x)))}
\end{Highlighting}
\end{Shaded}

\begin{verbatim}
##      [,1]
## [1,]    0
## [2,]    0
## [3,]    0
## [4,]    0
\end{verbatim}

\begin{Shaded}
\begin{Highlighting}[]
\KeywordTok{print}\NormalTok{(}\KeywordTok{zapsmall}\NormalTok{(}\KeywordTok{nullRC}\NormalTok{ (}\KeywordTok{t}\NormalTok{(x))))}
\end{Highlighting}
\end{Shaded}

\begin{verbatim}
##            [,1]
## [1,] -0.4467204
## [2,]  0.0655973
## [3,]  0.4633603
## [4,]  0.3315631
## [5,]  0.6866594
\end{verbatim}

\subsubsection{Singular}\label{singular-2}

\begin{Shaded}
\begin{Highlighting}[]
\KeywordTok{set.seed}\NormalTok{(}\DecValTok{12345}\NormalTok{)}
\NormalTok{x <-}\StringTok{ }\KeywordTok{matrix}\NormalTok{ (}\KeywordTok{rnorm}\NormalTok{(}\DecValTok{20}\NormalTok{), }\DecValTok{5}\NormalTok{, }\DecValTok{4}\NormalTok{)}
\NormalTok{x[,}\DecValTok{3}\NormalTok{] <-}\StringTok{ }\NormalTok{x[,}\DecValTok{1}\NormalTok{] }\OperatorTok{+}\StringTok{ }\NormalTok{x[,}\DecValTok{2}\NormalTok{]}
\NormalTok{x[,}\DecValTok{4}\NormalTok{] <-}\StringTok{ }\NormalTok{x[,}\DecValTok{1}\NormalTok{] }\OperatorTok{+}\StringTok{ }\NormalTok{x[,}\DecValTok{2}\NormalTok{]}
\KeywordTok{print}\NormalTok{(}\KeywordTok{zapsmall}\NormalTok{(}\KeywordTok{nullRC}\NormalTok{ (x)), }\DataTypeTok{digits =} \DecValTok{4}\NormalTok{)}
\end{Highlighting}
\end{Shaded}

\begin{verbatim}
##         [,1]    [,2]
## [1,] -0.5774 -0.2582
## [2,] -0.5774 -0.2582
## [3,]  0.0000  0.7746
## [4,]  0.5774 -0.5164
\end{verbatim}

\subsubsection{Broad}\label{broad-2}

\begin{Shaded}
\begin{Highlighting}[]
\KeywordTok{set.seed}\NormalTok{(}\DecValTok{12345}\NormalTok{)}
\NormalTok{x <-}\StringTok{ }\KeywordTok{matrix}\NormalTok{ (}\KeywordTok{rnorm}\NormalTok{(}\DecValTok{20}\NormalTok{), }\DecValTok{5}\NormalTok{, }\DecValTok{4}\NormalTok{)}
\KeywordTok{print}\NormalTok{(}\KeywordTok{zapsmall}\NormalTok{(}\KeywordTok{nullRC}\NormalTok{ (}\KeywordTok{t}\NormalTok{(x))), }\DataTypeTok{digits =} \DecValTok{4}\NormalTok{)}
\end{Highlighting}
\end{Shaded}

\begin{verbatim}
##         [,1]
## [1,] -0.4467
## [2,]  0.0656
## [3,]  0.4634
## [4,]  0.3316
## [5,]  0.6867
\end{verbatim}

\subsubsection{Weird}\label{weird-2}

\begin{Shaded}
\begin{Highlighting}[]
\NormalTok{x <-}\StringTok{ }\KeywordTok{matrix}\NormalTok{(}\DecValTok{1}\NormalTok{, }\DecValTok{1}\NormalTok{, }\DecValTok{6}\NormalTok{)}
\KeywordTok{print}\NormalTok{(}\KeywordTok{zapsmall}\NormalTok{(}\KeywordTok{nullRC}\NormalTok{ (x)), }\DataTypeTok{digits =} \DecValTok{4}\NormalTok{)}
\end{Highlighting}
\end{Shaded}

\begin{verbatim}
##         [,1]    [,2]    [,3]    [,4]    [,5]
## [1,] -0.7071 -0.4082 -0.2887 -0.2236 -0.1826
## [2,]  0.0000  0.0000  0.0000  0.0000  0.9129
## [3,]  0.7071 -0.4082 -0.2887 -0.2236 -0.1826
## [4,]  0.0000  0.8165 -0.2887 -0.2236 -0.1826
## [5,]  0.0000  0.0000  0.8660 -0.2236 -0.1826
## [6,]  0.0000  0.0000  0.0000  0.8944 -0.1826
\end{verbatim}

\begin{Shaded}
\begin{Highlighting}[]
\KeywordTok{print}\NormalTok{(}\KeywordTok{zapsmall}\NormalTok{(}\KeywordTok{nullRC}\NormalTok{ (}\KeywordTok{t}\NormalTok{(x))), }\DataTypeTok{digits =} \DecValTok{4}\NormalTok{)}
\end{Highlighting}
\end{Shaded}

\begin{verbatim}
##      [,1]
## [1,]    0
\end{verbatim}

\section{The Moore-Penrose Inverse}\label{the-moore-penrose-inverse}

\subsection{Theory}\label{theory-3}

The Moore-Penrose inverse of the \(n\times m\) matrix \(X\) is the
\(m\times n\) matrix \(X^+\), uniquely defined by the four Penrose
conditions

\begin{itemize}
\tightlist
\item
  \(XX^+\) is symmetric,
\item
  \(X^+X\) is symmetric,
\item
  \(XX^+X=X\)
\item
  \(X^+XX^+=X^+\).
\end{itemize}

We use MacDuffee's theorem, in combination with the full-rank version of
QR from theorem 1. A discussion of the origin of the lemma, and a proof,
are in Ben-Israel and Greville (2001), page 37.

\textbf{Theorem 4: {[}MacDuffee{]}} If \(X=FG\) is a full-rank
decomposition, then \(X^+=G'(F'XG')^{-1}F'\).

\textbf{Corollary 2: {[}General Inverse{]}} If \(X=QR\) is a full-rank
QR decomposition, then \(X^+=R'(RR')^{-1}Q'\).

\textbf{Proof:} Simply substitute \(X=QR\) in theorem 4. \textbf{QED}

Alternatively, from Cline (1964), using representation \(\eqref{E:qr2}\)
and \(U:=T^{-1}S\), \[
X^+=
P\begin{bmatrix}(I-U(I+U'U)^{-1}U')T^{-1}&0_{r\times(n-r)}\\(I+U'U)^{-1}U')T^{-1}&0_{(m-r)\times(n-r)}\end{bmatrix}\begin{bmatrix}Q'\\Q_\perp'\end{bmatrix}.
\]

\subsection{Computation}\label{computation-3}

The \texttt{ginv()} function in the package \texttt{MASS} (Venables and
Ripley (2002)) computes the Moore-Penrose inverse by using the singular
value decomposition. This is somewhat wasteful from the computational
point of view. In our function \texttt{ginvRC()} we use the result in
corollary 2.

\subsection{Examples}\label{examples-3}

\subsubsection{Regular}\label{regular-3}

\begin{Shaded}
\begin{Highlighting}[]
\KeywordTok{set.seed}\NormalTok{(}\DecValTok{12345}\NormalTok{)}
\NormalTok{x <-}\StringTok{ }\KeywordTok{matrix}\NormalTok{ (}\KeywordTok{rnorm}\NormalTok{(}\DecValTok{20}\NormalTok{), }\DecValTok{5}\NormalTok{, }\DecValTok{4}\NormalTok{)}
\KeywordTok{print}\NormalTok{(}\KeywordTok{ginvRC}\NormalTok{(x), }\DataTypeTok{digits =} \DecValTok{4}\NormalTok{)}
\end{Highlighting}
\end{Shaded}

\begin{verbatim}
##           [,1]   [,2]    [,3]     [,4]    [,5]
## [1,]  0.001437 0.5543 -1.1062 -0.08611  0.7360
## [2,] -0.475830 0.1896 -0.9106  0.17322  0.2032
## [3,]  0.152025 0.3173  0.2716  0.28814 -0.2538
## [4,] -0.058472 0.1057 -0.9417  0.66952  0.2640
\end{verbatim}

\subsubsection{Singular}\label{singular-3}

\begin{Shaded}
\begin{Highlighting}[]
\KeywordTok{set.seed}\NormalTok{(}\DecValTok{12345}\NormalTok{)}
\NormalTok{x <-}\StringTok{ }\KeywordTok{matrix}\NormalTok{ (}\KeywordTok{rnorm}\NormalTok{(}\DecValTok{20}\NormalTok{), }\DecValTok{5}\NormalTok{, }\DecValTok{4}\NormalTok{)}
\NormalTok{x[,}\DecValTok{3}\NormalTok{] <-}\StringTok{ }\NormalTok{x[,}\DecValTok{1}\NormalTok{] }\OperatorTok{+}\StringTok{ }\NormalTok{x[,}\DecValTok{2}\NormalTok{]}
\KeywordTok{print}\NormalTok{(}\KeywordTok{ginvRC}\NormalTok{(x), }\DataTypeTok{digits =} \DecValTok{4}\NormalTok{)}
\end{Highlighting}
\end{Shaded}

\begin{verbatim}
##          [,1]      [,2]    [,3]       [,4]    [,5]
## [1,]  0.21990  0.432249 -0.3261 -0.0008035  0.3222
## [2,] -0.29032 -0.001226 -0.1895  0.1960648 -0.1556
## [3,] -0.07043  0.431023 -0.5156  0.1952613  0.1666
## [4,] -0.01212  0.202422 -0.8589  0.7573697  0.1866
\end{verbatim}

\subsubsection{Broad}\label{broad-3}

\begin{Shaded}
\begin{Highlighting}[]
\KeywordTok{set.seed}\NormalTok{(}\DecValTok{12345}\NormalTok{)}
\NormalTok{x <-}\StringTok{ }\KeywordTok{matrix}\NormalTok{ (}\KeywordTok{rnorm}\NormalTok{(}\DecValTok{20}\NormalTok{), }\DecValTok{5}\NormalTok{, }\DecValTok{4}\NormalTok{)}
\KeywordTok{print}\NormalTok{(}\KeywordTok{ginvRC}\NormalTok{( }\KeywordTok{t}\NormalTok{(x)), }\DataTypeTok{digits =} \DecValTok{4}\NormalTok{)}
\end{Highlighting}
\end{Shaded}

\begin{verbatim}
##           [,1]    [,2]    [,3]     [,4]
## [1,]  0.001437 -0.4758  0.1520 -0.05847
## [2,]  0.554303  0.1896  0.3173  0.10568
## [3,] -1.106171 -0.9106  0.2716 -0.94165
## [4,] -0.086113  0.1732  0.2881  0.66952
## [5,]  0.736010  0.2032 -0.2538  0.26401
\end{verbatim}

\subsubsection{Weird}\label{weird-3}

\begin{Shaded}
\begin{Highlighting}[]
\NormalTok{x <-}\StringTok{ }\KeywordTok{matrix}\NormalTok{(}\DecValTok{1}\NormalTok{, }\DecValTok{1}\NormalTok{, }\DecValTok{6}\NormalTok{)}
\KeywordTok{print}\NormalTok{(}\KeywordTok{ginvRC}\NormalTok{(x), }\DataTypeTok{digits =} \DecValTok{4}\NormalTok{)}
\end{Highlighting}
\end{Shaded}

\begin{verbatim}
##        [,1]
## [1,] 0.1667
## [2,] 0.1667
## [3,] 0.1667
## [4,] 0.1667
## [5,] 0.1667
## [6,] 0.1667
\end{verbatim}

\section{Appendix: Code}\label{appendix-code}

\subsection{R code}\label{r-code}

\begin{Shaded}
\begin{Highlighting}[]
\KeywordTok{dyn.load}\NormalTok{(}\StringTok{"matrix.so"}\NormalTok{)}

\NormalTok{gsRC <-}\StringTok{ }\ControlFlowTok{function}\NormalTok{ (x, }\DataTypeTok{eps =} \FloatTok{1e-10}\NormalTok{) \{}
\NormalTok{  n <-}\StringTok{ }\KeywordTok{nrow}\NormalTok{ (x)}
\NormalTok{  m <-}\StringTok{ }\KeywordTok{ncol}\NormalTok{ (x)}
\NormalTok{  h <-}
\StringTok{    }\KeywordTok{.C}\NormalTok{(}
      \StringTok{"gsC"}\NormalTok{,}
      \DataTypeTok{x =} \KeywordTok{as.double}\NormalTok{(x),}
      \DataTypeTok{r =} \KeywordTok{as.double}\NormalTok{ (}\KeywordTok{matrix}\NormalTok{ (}\DecValTok{0}\NormalTok{, m, m)),}
      \DataTypeTok{n =} \KeywordTok{as.integer}\NormalTok{ (n),}
      \DataTypeTok{m =} \KeywordTok{as.integer}\NormalTok{ (m),}
      \DataTypeTok{rank =} \KeywordTok{as.integer}\NormalTok{ (}\DecValTok{0}\NormalTok{),}
      \DataTypeTok{pivot =} \KeywordTok{as.integer}\NormalTok{ (}\DecValTok{1}\OperatorTok{:}\NormalTok{m),}
      \DataTypeTok{eps =} \KeywordTok{as.double}\NormalTok{ (eps)}
\NormalTok{    )}
\NormalTok{  rank =}\StringTok{ }\NormalTok{h}\OperatorTok{$}\NormalTok{rank}
  \KeywordTok{return}\NormalTok{ (}\KeywordTok{list}\NormalTok{ (}
    \DataTypeTok{x =} \KeywordTok{matrix}\NormalTok{ (h}\OperatorTok{$}\NormalTok{x, n, m)[, }\DecValTok{1}\OperatorTok{:}\NormalTok{rank, }\DataTypeTok{drop =} \OtherTok{FALSE}\NormalTok{],}
    \DataTypeTok{r =} \KeywordTok{matrix}\NormalTok{ (h}\OperatorTok{$}\NormalTok{r, m, m)[}\DecValTok{1}\OperatorTok{:}\NormalTok{rank, , }\DataTypeTok{drop =} \OtherTok{FALSE}\NormalTok{],}
    \DataTypeTok{rank =}\NormalTok{ rank,}
    \DataTypeTok{pivot =}\NormalTok{ h}\OperatorTok{$}\NormalTok{pivot}
\NormalTok{  ))}
\NormalTok{\}}

\NormalTok{lsRC <-}\StringTok{ }\ControlFlowTok{function}\NormalTok{ (x, y, }\DataTypeTok{eps =} \FloatTok{1e-10}\NormalTok{) \{}
\NormalTok{  n <-}\StringTok{ }\KeywordTok{length}\NormalTok{ (y)}
\NormalTok{  m <-}\StringTok{ }\KeywordTok{ncol}\NormalTok{ (x)}
\NormalTok{  h <-}\StringTok{ }\KeywordTok{gsRC}\NormalTok{ (x, eps)}
\NormalTok{  l <-}\StringTok{ }\NormalTok{h}\OperatorTok{$}\NormalTok{rank}
\NormalTok{  p <-}\StringTok{ }\KeywordTok{order}\NormalTok{ (h}\OperatorTok{$}\NormalTok{pivot)}
\NormalTok{  k <-}\StringTok{ }\DecValTok{1}\OperatorTok{:}\NormalTok{l}
\NormalTok{  q <-}\StringTok{ }\NormalTok{h}\OperatorTok{$}\NormalTok{x}
\NormalTok{  a <-}\StringTok{ }\NormalTok{h}\OperatorTok{$}\NormalTok{r[, k, drop =}\StringTok{ }\OtherTok{FALSE}\NormalTok{]}
\NormalTok{  v <-}\StringTok{ }\NormalTok{h}\OperatorTok{$}\NormalTok{r[, }\OperatorTok{-}\NormalTok{k, drop =}\StringTok{ }\OtherTok{FALSE}\NormalTok{]}
\NormalTok{  u <-}\StringTok{ }\KeywordTok{crossprod}\NormalTok{ (q, y)}
\NormalTok{  b <-}\StringTok{ }\KeywordTok{solve}\NormalTok{ (a, u)}
\NormalTok{  res <-}\StringTok{ }\KeywordTok{drop}\NormalTok{ (y }\OperatorTok{-}\StringTok{ }\NormalTok{q }\OperatorTok{%*%}\StringTok{ }\NormalTok{u)}
\NormalTok{  s <-}\StringTok{ }\KeywordTok{sum}\NormalTok{ (res }\OperatorTok{^}\StringTok{ }\DecValTok{2}\NormalTok{)}
\NormalTok{  b <-}\StringTok{ }\KeywordTok{c}\NormalTok{(b, }\KeywordTok{rep}\NormalTok{ (}\DecValTok{0}\NormalTok{, m }\OperatorTok{-}\StringTok{ }\NormalTok{l))[p]}
  \ControlFlowTok{if}\NormalTok{ (l }\OperatorTok{==}\StringTok{ }\NormalTok{m) \{}
\NormalTok{    e <-}\StringTok{ }\KeywordTok{matrix}\NormalTok{(}\DecValTok{0}\NormalTok{, m, }\DecValTok{1}\NormalTok{)}
\NormalTok{  \} }\ControlFlowTok{else}\NormalTok{ \{}
\NormalTok{    e <-}\StringTok{ }\KeywordTok{rbind}\NormalTok{ (}\OperatorTok{-}\KeywordTok{solve}\NormalTok{(a, v), }\KeywordTok{diag}\NormalTok{(m }\OperatorTok{-}\StringTok{ }\NormalTok{l))[p, , drop =}\StringTok{ }\OtherTok{FALSE}\NormalTok{]}
\NormalTok{  \}}
  \KeywordTok{return}\NormalTok{ (}\KeywordTok{list}\NormalTok{ (}
    \DataTypeTok{solution =}\NormalTok{ b,}
    \DataTypeTok{residuals =}\NormalTok{ res,}
    \DataTypeTok{minssq =}\NormalTok{ s,}
    \DataTypeTok{nullspace =}\NormalTok{ e,}
    \DataTypeTok{rank =}\NormalTok{ l,}
    \DataTypeTok{pivot =}\NormalTok{ p}
\NormalTok{  ))}
\NormalTok{\}}

\NormalTok{nullRC <-}\StringTok{ }\ControlFlowTok{function}\NormalTok{ (x, }\DataTypeTok{eps =} \FloatTok{1e-10}\NormalTok{) \{}
\NormalTok{  h <-}\StringTok{ }\KeywordTok{gsRC}\NormalTok{ (x, }\DataTypeTok{eps =}\NormalTok{ eps)}
\NormalTok{  rank <-}\StringTok{ }\NormalTok{h}\OperatorTok{$}\NormalTok{rank}
\NormalTok{  p <-}\StringTok{ }\KeywordTok{order}\NormalTok{ (h}\OperatorTok{$}\NormalTok{pivot)}
\NormalTok{  r <-}\StringTok{ }\NormalTok{h}\OperatorTok{$}\NormalTok{r}
\NormalTok{  m <-}\StringTok{ }\KeywordTok{ncol}\NormalTok{ (x)}
\NormalTok{  t <-}\StringTok{ }\NormalTok{r[, }\DecValTok{1}\OperatorTok{:}\NormalTok{rank, drop =}\StringTok{ }\OtherTok{FALSE}\NormalTok{]}
\NormalTok{  s <-}\StringTok{ }\NormalTok{r[, }\OperatorTok{-}\NormalTok{(}\DecValTok{1}\OperatorTok{:}\NormalTok{rank), drop =}\StringTok{ }\OtherTok{FALSE}\NormalTok{]}
  \ControlFlowTok{if}\NormalTok{ (rank }\OperatorTok{==}\StringTok{ }\NormalTok{m)}
    \KeywordTok{return}\NormalTok{ (}\KeywordTok{matrix}\NormalTok{(}\DecValTok{0}\NormalTok{, m, }\DecValTok{1}\NormalTok{))}
  \ControlFlowTok{else}\NormalTok{ \{}
\NormalTok{    nullspace <-}\StringTok{ }\KeywordTok{rbind}\NormalTok{ (}\OperatorTok{-}\KeywordTok{solve}\NormalTok{(t, s), }\KeywordTok{diag}\NormalTok{ (m }\OperatorTok{-}\StringTok{ }\NormalTok{rank))[p, , drop =}\StringTok{ }\OtherTok{FALSE}\NormalTok{]}
    \KeywordTok{return}\NormalTok{ (}\KeywordTok{gsRC}\NormalTok{ (nullspace)}\OperatorTok{$}\NormalTok{x)}
\NormalTok{  \}}
\NormalTok{\}}

\NormalTok{ginvRC <-}\StringTok{ }\ControlFlowTok{function}\NormalTok{ (x, }\DataTypeTok{eps =} \FloatTok{1e-10}\NormalTok{) \{}
\NormalTok{  h <-}\StringTok{ }\KeywordTok{gsRC}\NormalTok{ (x, eps)}
\NormalTok{  p <-}\StringTok{ }\KeywordTok{order}\NormalTok{(h}\OperatorTok{$}\NormalTok{pivot)}
\NormalTok{  q <-}\StringTok{ }\NormalTok{h}\OperatorTok{$}\NormalTok{x}
\NormalTok{  s <-}\StringTok{ }\NormalTok{h}\OperatorTok{$}\NormalTok{r}
\NormalTok{  z <-}\StringTok{ }\KeywordTok{crossprod}\NormalTok{ (s, (}\KeywordTok{solve}\NormalTok{ (}\KeywordTok{tcrossprod}\NormalTok{(s), }\KeywordTok{t}\NormalTok{(q))))}
  \KeywordTok{return}\NormalTok{ (z[p, , }\DataTypeTok{drop =} \OtherTok{FALSE}\NormalTok{])}
\NormalTok{\}}
\end{Highlighting}
\end{Shaded}

\subsection{C code}\label{c-code}

\begin{Shaded}
\begin{Highlighting}[]
\PreprocessorTok{#include }\ImportTok{<math.h>}

\PreprocessorTok{#define MINDEX(i, j, n) (((j) - 1) * (n) + (i) - 1)}

\DataTypeTok{void}\NormalTok{ gsC(}\DataTypeTok{double}\NormalTok{ *, }\DataTypeTok{double}\NormalTok{ *, }\DataTypeTok{int}\NormalTok{ *, }\DataTypeTok{int}\NormalTok{ *, }\DataTypeTok{int}\NormalTok{ *, }\DataTypeTok{int}\NormalTok{ *, }\DataTypeTok{double}\NormalTok{ *);}

\DataTypeTok{void}\NormalTok{ gsC(}\DataTypeTok{double}\NormalTok{ *x, }\DataTypeTok{double}\NormalTok{ *r, }\DataTypeTok{int}\NormalTok{ *n, }\DataTypeTok{int}\NormalTok{ *m, }\DataTypeTok{int}\NormalTok{ *rank, }\DataTypeTok{int}\NormalTok{ *pivot,}
         \DataTypeTok{double}\NormalTok{ *eps) \{}
  \DataTypeTok{int}\NormalTok{ i, j, ip, nn = *n, mm = *m, rk = *m, jwork = }\DecValTok{1}\NormalTok{;}
  \DataTypeTok{double}\NormalTok{ s = }\FloatTok{0.0}\NormalTok{, p;}
  \ControlFlowTok{for}\NormalTok{ (j = }\DecValTok{1}\NormalTok{; j <= mm; j++) \{}
\NormalTok{    pivot[j - }\DecValTok{1}\NormalTok{] = j;}
\NormalTok{  \}}
  \ControlFlowTok{while}\NormalTok{ (jwork <= rk) \{}
    \ControlFlowTok{for}\NormalTok{ (j = }\DecValTok{1}\NormalTok{; j < jwork; j++) \{}
\NormalTok{      s = }\FloatTok{0.0}\NormalTok{;}
      \ControlFlowTok{for}\NormalTok{ (i = }\DecValTok{1}\NormalTok{; i <= nn; i++) \{}
\NormalTok{        s += x[MINDEX(i, jwork, nn)] * x[MINDEX(i, j, nn)];}
\NormalTok{      \}}
\NormalTok{      r[MINDEX(j, jwork, mm)] = s;}
      \ControlFlowTok{for}\NormalTok{ (i = }\DecValTok{1}\NormalTok{; i <= nn; i++) \{}
\NormalTok{        x[MINDEX(i, jwork, nn)] -= s * x[MINDEX(i, j, nn)];}
\NormalTok{      \}}
\NormalTok{    \}}
\NormalTok{    s = }\FloatTok{0.0}\NormalTok{;}
    \ControlFlowTok{for}\NormalTok{ (i = }\DecValTok{1}\NormalTok{; i <= nn; i++) \{}
\NormalTok{      s += x[MINDEX(i, jwork, nn)] * x[MINDEX(i, jwork, nn)];}
\NormalTok{    \}}
    \ControlFlowTok{if}\NormalTok{ (s > *eps) \{}
\NormalTok{      s = sqrt(s);}
\NormalTok{      r[MINDEX(jwork, jwork, mm)] = s;}
      \ControlFlowTok{for}\NormalTok{ (i = }\DecValTok{1}\NormalTok{; i <= nn; i++) \{}
\NormalTok{        x[MINDEX(i, jwork, nn)] /= s;}
\NormalTok{      \}}
\NormalTok{      jwork += }\DecValTok{1}\NormalTok{;}
\NormalTok{    \}}
    \ControlFlowTok{if}\NormalTok{ (s <= *eps) \{}
\NormalTok{      ip = pivot [rk - }\DecValTok{1}\NormalTok{];}
\NormalTok{      pivot[rk - }\DecValTok{1}\NormalTok{] = pivot[jwork - }\DecValTok{1}\NormalTok{];}
\NormalTok{      pivot[jwork - }\DecValTok{1}\NormalTok{] = ip;}
      \ControlFlowTok{for}\NormalTok{ (i = }\DecValTok{1}\NormalTok{; i <= nn; i++) \{}
\NormalTok{        p = x[MINDEX(i, rk, nn)];}
\NormalTok{        x[MINDEX(i, rk, nn)] = x[MINDEX(i, jwork, nn)];}
\NormalTok{        x[MINDEX(i, jwork, nn)] = p;}
\NormalTok{      \}}
      \ControlFlowTok{for}\NormalTok{ (j = }\DecValTok{1}\NormalTok{; j <= mm; j++) \{}
\NormalTok{        p = r[MINDEX(j, rk, mm)];}
\NormalTok{        r[MINDEX(j, rk, mm)] = r[MINDEX(j, jwork, mm)];}
\NormalTok{        r[MINDEX(j, jwork, mm)] = p;}
\NormalTok{      \}}
\NormalTok{      rk -= }\DecValTok{1}\NormalTok{;}
\NormalTok{    \}}
\NormalTok{  \}}
\NormalTok{  *rank = rk;}
\NormalTok{\}}
\end{Highlighting}
\end{Shaded}

\section{NEWS}\label{news}

001 02/28/16

\begin{itemize}
\tightlist
\item
  First release
\end{itemize}

002 03/20/16

\begin{itemize}
\tightlist
\item
  Corrections and clarifications suggested by Bill Venables
\end{itemize}

003 03/21/16

\begin{itemize}
\tightlist
\item
  Example with explicit permutation (to set my mind at ease)
\item
  Computed both left and right null spaces
\end{itemize}

\section*{References}\label{references}
\addcontentsline{toc}{section}{References}

\hypertarget{refs}{}
\hypertarget{ref-benisrael_greville_01}{}
Ben-Israel, A., and T.N.E. Greville. 2001. \emph{Generalized Inverses:
Theory and Applications}. Second Edition. New York: Springer.

\hypertarget{ref-bjorck_96}{}
Björck, Å. 1996. \emph{Numerical Methods for Least Squares Problems}.
Philadelphia, PA: SIAM.

\hypertarget{ref-cline_64}{}
Cline, R.E. 1964. ``Representations for the Generalized Inverse of a
Partitioned Matrix.'' \emph{Journal of the Society for Industrial and
Applied Mathematics} 12 (3): 588--600.

\hypertarget{ref-deleeuw_E_15c}{}
De Leeuw, J. 2015. ``Exceedingly Simple Gram-Schmidt Code.''
\url{http://deleeuwpdx.net/pubfolders/gs/gs.pdf}.

\hypertarget{ref-venables_ripley_02}{}
Venables, W.N., and B.D. Ripley. 2002. \emph{Modern Applied Statistics
with S}. Fourth Edition. New York: Springer.
\href{\%7Bhttp://www.stats.ox.ac.uk/pub/MASS4\%7D}{\{http://www.stats.ox.ac.uk/pub/MASS4\}}.


\end{document}
